% ============================================================
% 04_experimental_setup.tex  (revised 2026-10)
% ============================================================

% ------------------------------------------------------------
\subsection{Yulara: a curtailed off-grid mini-grid}
\label{subsec:yulara}

The Yulara installation (Northern Territory, Australia; $25.24^\circ$\,S, $131.04^\circ$\,E,
492~m, hot desert climate) consists of five PV sites (1.8~MW in total) feeding the
gas-supplied mini-grid of the Ayers Rock Resort; per-site 5-min measurements and a shared
weather station are published by the DKA Solar Centre~\cite{dkasc_yulara}. The series used in
earlier work on this plant, including the first version of the present study, is exactly the
output of the largest site, Desert Gardens (1{,}058~kW). A site-by-site analysis shows that this
site is systematically curtailed around solar noon (Fig.~\ref{fig:curtailment}): under high
irradiance its output relative to the other sites falls to 0.37 at noon while the four other
sites stay between 0.98 and 1.07, and on clear days the total PV output is held at
615--630~kW between 11:00 and 15:00. Treating curtailment as censoring of the available power
(delivered $=\min$(available, cap)) and reconstructing the available power from the
uncurtailed sites, we estimate that 50.6\% of the energy available at Desert Gardens was
curtailed (conservative bounds 31.7--59.1\%; 36--39\% in summer, 56--65\% from April to
September). These figures are inferred from the data; the control logic has not been confirmed
by the operator.

A curtailed series mixes weather and dispatch decisions: its $k_t$ saturates at the
curtailment cap instead of the clear-sky envelope, and a forecaster trained on it learns the
dispatch rule. The forecasting target is therefore the summed output of the four uncurtailed
sites (Service Station 227~kW, Sails in the Desert 107~kW, Laundry 328~kW and Connellan Airport
106~kW; 768~kW in total) from 2 April 2016 to 13 May 2020, after which these sites stop
reporting for six months while Desert Gardens keeps operating.

\begin{figure}[t]
  \centering
  \includegraphics[width=0.85\linewidth]{fig_curtailment.pdf}
  \caption{Output of each Yulara site relative to the others under high irradiance, by hour of
  the day (normalised to 08--09~h). Desert Gardens, the series used in earlier studies, falls to
  0.37 of its expected share at noon; Connellan Airport's monotone afternoon decline is an
  orientation effect.}
  \label{fig:curtailment}
\end{figure}

% ------------------------------------------------------------
\subsection{NIST: a grid-connected array in a humid subtropical climate}
\label{subsec:nist}

The second site is the Ground array of the NIST campus in Gaithersburg, Maryland, USA
($39.13^\circ$\,N, $77.21^\circ$\,W, 138~m; K\"oppen Cfa), 271~kW$_\mathrm{DC}$, south-facing,
tilted at $20^\circ$~\cite{boyd2017nist}. Its 1-min inverter AC power and the array's own
sensors (horizontal thermopile pyranometer, ambient temperature, wind) were obtained from the
public PVDAQ data lake (system 4902, NIST\_Ground\_1)~\cite{deline2021pvdaq}, whose values are
identical to the NIST archive for the overlapping year 2016, except the pyranometer, stored in
millivolts. Its conversion to W\,m$^{-2}$ follows NIST's own constants (one change in May 2016,
reproduced exactly from the NIST archive); the raw signal drops by about 16\% between 10 and
17~October~2017 while the calibrated reference cell stays consistent with the AC power, which
indicates a sensor change with an unknown constant. The series is therefore restricted to
1~January~2016 -- 10~October~2017. An inverter outage (21~June -- 5~July~2016, 524~h) is
excluded by the gap rule. Pressure is not measured at the array and is set to the standard
pressure of the altitude (a constant feature).

% ------------------------------------------------------------
\subsection{Weather forecasts}
\label{subsec:gfs_data}

GFS forecasts were retrieved from the NSF NCAR Geoscience Data Exchange (dataset
d084001~\cite{ncep2015gfs}) as monthly subsets of all runs with forecast hours up to 48~h, for
April~2016 -- May~2020 at Yulara and January~2016 -- December~2018 at NIST. Every target of
every test forecast has a GFS value from the run selected by Eq.~\eqref{eq:run}; 0.1\% of the
NIST training and validation targets have none (the tree models handle these missing values
natively).

% ------------------------------------------------------------
\subsection{Partitions, runs and reproducibility}
\label{subsec:partitions}

Each series is split chronologically into training, validation and test partitions (60/20/20\%
of the rows; Table~\ref{tab:partitions}). The test partition contains only forecasts issued after
the end of the validation period. Every configuration is run with and without GFS; the neural
members of the run with GFS are the checkpoints of the run without GFS, so that the two runs
differ only by the GFS information. At Yulara, five seeds (42, 1, 2, 3, 4) control the neural
networks and the tree models. As a sensitivity analysis, a season-balanced validation keeps the
same test partition but uses every fourth week of the preceding period as validation, dropping
sequences whose window mixes training and validation rows. Data preparation, models, selection,
evaluation and every table of this paper are produced by the scripts and configuration files of
the public repository.

\begin{table}[t]
  \centering
  \caption{Data partitions (first and last target day of the forecasts in each partition;
  number of hourly forecasts issued, each with 24 lead times).}
  \label{tab:partitions}
  \small
  \begin{tabular}{lllll}
    \toprule
    Site & Capacity & Training & Validation & Test \\
    \midrule
    Yulara (4 sites) & 768 kW & 2016-04-03 -- 2018-09-19 & 2018-09-21 -- 2019-07-17 & 2019-07-19 -- 2020-05-13 \\
                     &        & 19{,}847 & 5{,}699 & 6{,}650 \\
    NIST Ground      & 271 kW & 2016-01-02 -- 2017-01-23 & 2017-01-25 -- 2017-06-02 & 2017-06-04 -- 2017-10-10 \\
                     &        & 8{,}155 & 3{,}040 & 3{,}069 \\
    \bottomrule
  \end{tabular}
\end{table}

% ------------------------------------------------------------
\subsection{Evaluation}
\label{subsec:evaluation}

Forecasts are evaluated in power, $\hat P = \hat k_t P_\mathrm{cs}$, against the measured
power, on daytime targets ($P_\mathrm{cs}$ above the night threshold; measured zero output by day
is negligible, 2 of 15{,}671 hours above $10^\circ$ elevation at Yulara). We report the MAE and
RMSE in kW and as a percentage of the DC capacity (nMAE, nRMSE), and the skill score
$1 - \mathrm{MSE}/\mathrm{MSE}_\mathrm{ref}$ against day-ahead persistence. The coefficient of
determination of the power is high for every model because the diurnal cycle dominates it
(0.84 for day-ahead persistence at Yulara) and discriminates little between models; MAE in
$k_t$ units is given in the supplementary tables only, because it weighs dawn and dusk hours,
where $k_t$ is noisy and the power small, as much as midday hours and can reverse the ranking of
close models.

Forecasts issued on the same day share most of their targets, so losses are first averaged per
issue day. Differences between two forecasts are assessed with a bootstrap over days (95\%
interval of the MAE difference, 5{,}000 replicates) and a Diebold--Mariano
test~\cite{diebold1995comparing} on the daily loss differentials with a Newey--West
variance~\cite{newey1987simple}.
